\SourcePage{727}{730}
\section{The Breit--Wigner formulas}

In \S~134 quasistationary states were introduced as states having a finite
but comparatively long lifetime. A broad class of such states occurs in
nuclear reactions at moderate energies which proceed through the formation
of a \emph{compound nucleus}\footnote{The concept of the compound nucleus was
introduced by N.~Bohr (1936).}.

The physical picture is that the incident particle interacts with the
nucleons and ``merges'' with the nucleus, forming a compound system in which
its energy is distributed among many nucleons. Resonance energies correspond
to quasidiscrete levels of this system. Their lifetimes are long compared
with nuclear periods because, most of the time, the energy is shared among so
many particles that none can overcome the attraction of the others and leave
the nucleus. Only rarely is enough energy concentrated on one particle. The
compound nucleus can decay in various ways, corresponding to the possible
reaction channels\footnote{The competing processes also include radiative
capture of the incident particle, in which the excited compound nucleus emits
a $\gamma$ quantum and passes to its ground state. This process too is
``slow,'' because a radiative transition has comparatively low probability.}.

\SourcePage{728}{731}
This character of the collision shows that inelastic processes do not affect
the potential part of the elastic amplitude, which is unrelated to compound-
nucleus properties (see \S~134); they change only its resonance part. For the
same reason, inelastic amplitudes involving a compound nucleus are purely
resonant. The resonance denominator retains the form
$E-E_0+i\Gamma/2$, associated with the vanishing of the incoming-wave
coefficient at $E=E_0-i\Gamma/2$, and $\Gamma$ still determines the total
decay probability of the quasistationary state by any route.

These considerations, together with unitarity, determine the amplitudes.
Number all decay channels symmetrically, without first choosing the entrance
channel, using indices $a,b,c,\ldots$. Consider partial amplitudes for the
angular momentum $l$ of the quasistationary state\footnote{We first disregard
the complicating influence of particle spins.}. Seek them in the form
\begin{equation}
 f_{ab}^{(l)}=\frac1{2ik_a}(e^{2i\delta_a}-1)\delta_{ab}
 -\frac1{2\sqrt{k_ak_b}}e^{i(\delta_a+\delta_b)}
 \frac{\Gamma M_{ab}}{E-E_0+(1/2)i\Gamma}.
 \tag{145.1}
\end{equation}
The index $l$ on $\delta_a$ and $M_{ab}$ is suppressed. The first term occurs
only for $a=b$ and is the potential elastic amplitude in channel $a$; the
$\delta_a$ are the phases $\delta_l^{(0)}$ of (134.12). The second term is
resonant. Its coefficient is written in this form to simplify unitarity.

Since the fixed magnitude of orbital angular momentum is unchanged by time
reversal, reciprocity here simply requires $f_{ab}^{(l)}$ to be symmetric in
$a,b$. Thus $M_{ab}=M_{ba}$.

\SourcePage{729}{732}
The unitarity conditions are
\begin{equation}
 \operatorname{Im}f_{ab}^{(l)}=\sum_c k_c f_{ac}^{(l)}f_{bc}^{(l)*}
 \tag{145.2}
\end{equation}
(compare (144.8)). Substitution of (145.1) gives
\[
 \frac{M_{ab}^*}{E-E_0-(1/2)i\Gamma}
 -\frac{M_{ab}}{E-E_0+(1/2)i\Gamma}
 =\frac{i\Gamma\sum_c M_{ac}M_{bc}^*}
 {(E-E_0)^2+(1/4)\Gamma^2}.
\]
For this to hold identically in $E$, first $M_{ab}=M_{ab}^*$, so the
$M_{ab}$ are real, and then
\begin{equation}
 M_{ab}=\sum_c M_{ac}M_{bc}.
 \tag{145.3}
\end{equation}
Thus the matrix $M$ equals its square.

A real symmetric $M$ can be diagonalized by an orthogonal transformation
$U$. Denoting its eigenvalues by $M^{(\alpha)}$, write
\[
 \sum_{a,b}U_{\alpha a}U_{\beta b}M_{ab}=M^{(\alpha)}\delta_{\alpha\beta},
\]
with
\begin{equation}
 \sum_c U_{\alpha c}U_{\beta c}=\delta_{\alpha\beta}.
 \tag{145.4}
\end{equation}
Conversely,
\begin{equation}
 M_{ab}=\sum_\alpha U_{\alpha a}U_{\alpha b}M^{(\alpha)}.
 \tag{145.5}
\end{equation}
Equation (145.3) gives $M^{(\alpha)}=(M^{(\alpha)})^2$, so every eigenvalue is
0 or 1. If only $M^{(1)}=1$ is nonzero, then
\begin{equation}
 M_{ab}=U_{1a}U_{1b}.
 \tag{145.6}
\end{equation}

\SourcePage{730}{733}
If several eigenvalues are nonzero, $M_{ab}$ is a sum of terms involving
independent sets $U_{1a},U_{2a},\ldots$, constrained only by orthogonality.
This would describe accidental degeneracy, with several compound-nucleus
quasistationary states at the same quasidiscrete energy\footnote{This is
especially clear if every $M^{(\alpha)}=1$. Equations (145.4), (145.5) then
give $M_{ab}=\delta_{ab}$, so there are no transitions between different
channels. The case would represent several independent quasidiscrete states,
each realized in elastic scattering in one channel.}. Excluding these
uninteresting degenerate levels, the elements of $M$ factor into quantities
depending on one channel each.

Define
\[
 |U_{1a}|=\sqrt{\Gamma_a/\Gamma}.
\]
Then (145.6) becomes
\begin{equation}
 M_{ab}=\pm\sqrt{\Gamma_a\Gamma_b}/\Gamma,
 \tag{145.7}
\end{equation}
where the sign depends on those of $U_{1a}$ and $U_{1b}$. Since
$\sum_cU_{1c}U_{1c}=1$,
\begin{equation}
 \sum_a\Gamma_a=\Gamma.
 \tag{145.8}
\end{equation}
The $\Gamma_a$ are the \emph{partial widths} of the channels. Equations
(145.1), (145.7), (145.8) give the required general amplitudes.

Now select one entrance channel\footnote{These formulas were first obtained
by \emph{Breit and Wigner} (O.~Breit, E.~Wigner, 1936).}. Denote its partial
width by $\Gamma_e$, the \emph{elastic width}, and the widths of the reactions
by $\Gamma_{r_1},\Gamma_{r_2},\ldots$. The total elastic amplitude is
\begin{equation}
 f_e(\theta)=f^{(0)}(\theta)-\frac{2l+1}{2k}
 \frac{\Gamma_e}{E-E_0+(1/2)i\Gamma}e^{2i\delta_l^{(0)}}P_l(\cos\theta),
 \tag{145.9}
\end{equation}

\SourcePage{731}{734}
where $\mathbf k$ is the incident wave vector and $f^{(0)}$ is the potential
scattering amplitude. This differs from (134.12) by replacing $\Gamma$ in the
resonance numerator with the smaller $\Gamma_e$.

The inelastic amplitudes are purely resonant. Their differential and integral
cross-sections are
\begin{equation}
 d\sigma_{r_a}=\frac{(2l+1)^2}{4k^2}
 \frac{\Gamma_e\Gamma_{r_a}}{(E-E_0)^2+(1/4)\Gamma^2}
 [P_l(\cos\theta)]^2d\Omega,
 \tag{145.10}
\end{equation}
\begin{equation}
 \sigma_{r_a}=(2l+1)\frac\pi{k^2}
 \frac{\Gamma_e\Gamma_{r_a}}{(E-E_0)^2+(1/4)\Gamma^2}.
 \tag{145.11}
\end{equation}
The total inelastic cross-section is
\begin{equation}
 \sigma_r=(2l+1)\frac\pi{k^2}
 \frac{\Gamma_e\Gamma_r}{(E-E_0)^2+(1/4)\Gamma^2},
 \tag{145.12}
\end{equation}
where $\Gamma_r=\Gamma-\Gamma_e$ is the total inelastic width.

Integrating over the resonance region, and extending $E-E_0$ from
$-\infty$ to $\infty$ because $\sigma_r$ falls rapidly, gives
\begin{equation}
 \int\sigma_r\,dE=(2l+1)\frac{2\pi^2}{k^2}\frac{\Gamma_e\Gamma_r}{\Gamma}.
 \tag{145.13}
\end{equation}

For slow neutrons only $s$-wave scattering matters, and the potential
amplitude is the real constant $-\alpha$. Instead of (134.14),
\begin{equation}
 f_e=-\alpha-\frac{\Gamma_e}{2k(E-E_0+(1/2)i\Gamma)}.
 \tag{145.14}
\end{equation}
The total elastic cross-section is
\begin{equation}
 \sigma_e=4\pi\alpha^2+\frac\pi{k^2}
 \frac{\Gamma_e^2+4\alpha k\Gamma_e(E-E_0)}{(E-E_0)^2+(1/4)\Gamma^2}.
 \tag{145.15}
\end{equation}
The first term is the potential-scattering cross-section. Potential and
resonance scattering interfere. Only very near the level,
$E-E_0\sim\Gamma$, may $\alpha$ be negligible (recall
$|\alpha k|\ll1$), giving
\begin{equation}
 \sigma_e=\frac\pi{k^2}\frac{\Gamma_e^2}{(E-E_0)^2+(1/4)\Gamma^2}.
 \tag{145.16}
\end{equation}

\SourcePage{732}{735}
The total elastic plus inelastic cross-section is then
\begin{equation}
 \sigma_t=\sigma_e+\sigma_r=\frac\pi{k^2}
 \frac{\Gamma_e\Gamma}{(E-E_0)^2+(1/4)\Gamma^2}.
 \tag{145.17}
\end{equation}
When potential scattering is negligible,
\[
 \sigma_e=\sigma_t\frac{\Gamma_e}{\Gamma},\qquad
 \sigma_{r_a}=\sigma_t\frac{\Gamma_{r_a}}{\Gamma}.
\]
Here $\sigma_t$, the sum over every resonant process, is the cross-section for
forming the compound nucleus. Each elastic or inelastic cross-section is
$\sigma_t$ times the relative probability for the corresponding decay,
given by the partial width divided by the total width. This factorization of
$M_{ab}$ matches the two-stage physical picture: formation of a compound
nucleus in a quasistationary state, followed by decay through one
channel\footnote{The calculations above referred to reactions
$a+X=b+Y$, with two particles both before and after. This assumption is not
fundamental, as is clear from the physical result. Formulas such as (145.11)
for integral cross-sections also apply when more than one particle is emitted.}.

As noted in \S~134, these formulas require only
$|E-E_0|\ll D$, where $D$ is the spacing of neighboring compound-nucleus
levels with the same angular momentum. In this form, however, they do not
permit the limit $E\to0$ when zero lies in the resonance region. One must then
replace $E_0$ by a related constant $\varepsilon_0$, and
$\Gamma_e$ by $\gamma_e\sqrt E$, while keeping $\Gamma_r$ constant
(H.~A.~Bethe, G.~Placzek, 1937)\footnote{For inelastic processes possible at
low energies, such as radiative capture, $E=0$ is not a threshold. A similar
replacement for $\Gamma_{r_a}$ would be required near the threshold of that
reaction, below which it is impossible.}.

\SourcePage{733}{736}
The resulting inelastic cross-section (145.12) grows as $1/\sqrt E$ when
$E\to0$, in agreement with \S~143.

Spin generally gives cumbersome formulas. Consider only slow-neutron
scattering, where $l=0$. Coupling the target spin $i$ to the neutron spin
$s=1/2$ gives compound-nucleus spin $j=i\pm1/2$; assume $i\ne0$. Each
quasidiscrete level has one definite $j$. The reaction cross-section is
therefore (145.12), with $l=0$, multiplied by the probability $g(j)$ that the
nucleus--neutron system has the required $j$.

Suppose neutron and target spins are randomly oriented. There are
$2(2i+1)$ orientations of the pair, and $2j+1$ correspond to a specified
$j$. Hence
\begin{equation}
 g(j)=\frac{2j+1}{2(2i+1)}.
 \tag{145.18}
\end{equation}
The elastic formula changes similarly. Both $j$ values contribute to
potential scattering, so multiply the second term of (145.15) by the
$g(j)$ of the resonant level, and replace $4\pi\alpha^2$ by
$\sum_jg(j)\,4\pi\alpha^{(j)2}$.

Compound-nucleus formation also implies general properties of the angular
distribution of reaction products. Each quasistationary state has a definite
parity, which is inherited by the final system $(b+Y)$. Under inversion its
wave function and reaction amplitude can only acquire a factor $\pm1$, so
the cross-section is unchanged. In the center-of-mass system inversion sends

\SourcePage{734}{737}
$\theta\to\pi-\theta$, $\varphi\to\pi+\varphi$. Thus the distribution is
invariant under this replacement. After averaging over all spin directions,
it depends only on $\theta$ and is symmetric under
$\theta\to\pi-\theta$\footnote{For spinless particles the differential
reaction cross-section is simply proportional to
$[P_l(\cos\theta)]^2$, making the symmetry evident.}.

The great number of densely spaced compound-nucleus levels makes the detailed
energy dependence of scattering cross-sections extremely complicated and
obscures systematic changes between nuclei. It is therefore useful to average
over energy intervals large compared with level spacings, suppressing the
resonance detail. We also cease to distinguish types of inelastic process and
divide scattering, in the sense described below, only into ``elastic'' and
``inelastic'' parts\footnote{The averaging method leading to the optical model
of nuclear scattering was proposed by \emph{Weisskopf, Porter, and Feshbach}
(V.~F.~Weisskopf, C.~E.~Porter, H.~Feshbach, 1954).}.

To explain the averaging, again disregard spin and consider $l=0$. From
(142.7),
\begin{equation}
 \sigma_e=\frac\pi{k^2}|S-1|^2,\qquad
 \sigma_r=\frac\pi{k^2}(1-|S|^2),\qquad
 \sigma_t=\frac\pi{k^2}\,2(1-\operatorname{Re}S),
 \tag{145.19}
\end{equation}
where the superscript $(0)$ is suppressed. Since $\sigma_t$ is linear in $S$,
its energy average is
\begin{equation}
 \overline{\sigma}_t=\frac\pi{k^2}\,2(1-\operatorname{Re}\overline S),
 \tag{145.20}
\end{equation}
leaving the slowly varying $k^{-2}$ outside the average. Define the
``elastic'' cross-section in the averaged picture as

\SourcePage{735}{738}
\begin{equation}
 \overline{\sigma}_e^{\mathrm{opt}}=(\pi/k^2)|\overline S-1|^2,
 \tag{145.21}
\end{equation}
which generally differs from $\overline{\sigma_e}$. Thus the outgoing-wave
amplitude $Se^{ikr}/r$ is averaged before defining elastic scattering. This
preserves the form of an elastically scattered wave packet and describes the
``coherent'' part. Elastic scattering through the long-lived compound nucleus
is excluded, since its formation and decay erase the incident packet's
specific character. Define the averaged ``inelastic'' cross-section by
$\overline{\sigma}_a^{\mathrm{opt}}
=\overline{\sigma}_t-\overline{\sigma}_e^{\mathrm{opt}}$:
\begin{equation}
 \overline{\sigma}_a^{\mathrm{opt}}=(\pi/k^2)(1-|\overline S|^2).
 \tag{145.22}
\end{equation}
It includes both truly inelastic processes and elastic scattering through an
intermediate compound nucleus.

This interpretation gives the correct limiting pictures and provides a
reasonable interpolation. For well-resolved low-energy resonances
($\Gamma\ll D$), near each level
\[
 S=\left(1-\frac{i\Gamma_e}{E-E_0+(1/2)i\Gamma}\right)
 \exp(2i\delta^{(0)}).
\]
Averaging gives
\begin{equation}
 \overline S=(1-\pi\overline\Gamma_e/D)\exp(2i\delta^{(0)}),
 \tag{145.23}
\end{equation}
where $\overline\Gamma_e$ and $D$ are the mean elastic width and mean level
spacing in the interval, and slowly varying $\delta^{(0)}(E)$ is held
constant. Hence
\begin{equation}
 \overline{\sigma}_a^{\mathrm{opt}}=(\pi/k^2)(2\pi\overline\Gamma_e/D),
 \tag{145.24}
\end{equation}
neglecting terms of order $\Gamma/D$\footnote{Terms from the influence of
other levels near a given level would be of the same order.}. This equals the
mean of (145.17), the compound-nucleus formation cross-section.

\SourcePage{736}{739}
As excitation energy increases, level spacings decrease and decay
probabilities, hence total widths, increase, until levels overlap and the
notion of quasidiscrete levels largely loses meaning. Irregularities in
$S(E)$ are smoothed, so the exact and averaged functions become close, and
(145.22) coincides with $\sigma_r$ in (145.19). At high energies, decay back
through the entrance channel is negligible beside the many other decay modes;
all compound-nucleus processes may then be regarded as inelastic.

Thus the averaged scattering is again determined by one quantity,
$\overline S$, now a smooth function of energy. In the \emph{optical model},
the nucleus is approximated by a force field with a complex potential. Its
imaginary part produces absorption in addition to elastic scattering. This
absorption, with cross-section (145.22), is identified with ``inelastic''
scattering in the averaged picture.
